\setcounter{section}{0}

\Needspace{5\baselineskip}
\hypertarget{imitation-learning}{%
\section{Imitation Learning}\label{imitation-learning}}

\Needspace{5\baselineskip}
\hypertarget{i.-behavioral-cloning-bc}{%
\subsection{I. Behavioral Cloning (BC)}\label{i.-behavioral-cloning-bc}}

Behavioral cloning is the most intuitive and simplest form of imitation learning. Its core idea is to turn an imitation learning problem into a supervised learning problem.

\Needspace{5\baselineskip}
\hypertarget{i-algorithmic-principles}{%
\subsubsection{(I) Algorithmic Principles}\label{i-algorithmic-principles}}

\begin{enumerate}
\def\labelenumi{\arabic{enumi}.}
\item
  Core idea: behavioral cloning assumes that expert data are independent and identically distributed samples. It does not concern itself with the environment's dynamics model and directly establishes a mapping from state space \(S\) to action space \(A\).
\item
  Inputs and labels: we use states \(s_t\) in expert trajectories as sample inputs and the expert's actions \(a_t\) in those states as labels.
\item
\Needspace{5\baselineskip}
  Optimization objective: train a policy network whose outputs are as close as possible to the expert's actions. Mathematically, this is equivalent to minimizing the expected loss on expert dataset \(\mathcal{B}\):
\end{enumerate}

\[
\theta^{*}
=
\arg\min_\theta
\mathbb{E}_{(s,a)\sim\mathcal{B}}
\bigl[
L(\pi_\theta(s),a)
\bigr]
\]

\Needspace{5\baselineskip}
\hypertarget{ii-application-mainly-used-early-in-training-to-stabilize-model-training-more-quickly.-for-example-early-alphago-first-performed-imitation-learning-on-human-masters-game-records-then-used-other-methods-to-explore-new-policies-autonomously.}{%
\subsubsection{(II) Application: Mainly Used Early in Training to Stabilize Model Training More Quickly. For Example, Early AlphaGo First Performed Imitation Learning on Human Masters' Game Records, Then Used Other Methods to Explore New Policies Autonomously.}\label{ii-application-mainly-used-early-in-training-to-stabilize-model-training-more-quickly.-for-example-early-alphago-first-performed-imitation-learning-on-human-masters-game-records-then-used-other-methods-to-explore-new-policies-autonomously.}}

\Needspace{5\baselineskip}
\hypertarget{iii-core-limitation-accumulating-errors}{%
\subsubsection{(III) Core Limitation: Accumulating Errors}\label{iii-core-limitation-accumulating-errors}}

During training, the policy network sees only the state distribution generated by experts. During testing (actual inference), however, it inevitably makes small prediction errors. Once the agent acts according to an imperfect policy, it enters a new state. This state may deviate slightly from the expert trajectory distribution. Because the agent has never seen this ``state after deviation'' in the training set (for example, a different board position that looks similar to the original but is actually very different), it cannot correct the error and may instead take an even more inappropriate action. Such errors grow in a cascade over time steps, causing the agent's trajectory to diverge progressively from the expert's.

\Needspace{5\baselineskip}
\hypertarget{ii.-generative-adversarial-imitation-learning-gail-standford-2016}{%
\subsection{II. Generative Adversarial Imitation Learning (GAIL) (Standford, 2016)}\label{ii.-generative-adversarial-imitation-learning-gail-standford-2016}}

\Needspace{5\baselineskip}
\hypertarget{i-core-principle-occupancy-measure-matching}{%
\subsubsection{(I) Core Principle: Occupancy Measure Matching}\label{i-core-principle-occupancy-measure-matching}}

GAIL holds that if a policy is sufficiently good, the state-action occupancy measure \(\rho(s,a)\) it generates in the environment should match the expert policy's occupancy measure \(\rho_E(s,a)\).

\begin{itemize}
\tightlist
\item
  Occupancy measure \(\rho(s,a)\): it can be understood as the frequency distribution of particular state-action pairs over a long run.
\item
  Essence: GAIL uses adversarial training to reduce the distance between the distribution \(\rho_\pi\) generated by the agent's policy and the expert distribution \(\rho_E\), usually using Jensen-Shannon divergence.
\end{itemize}

\Needspace{5\baselineskip}
\hypertarget{ii-architecture-and-mathematical-derivation}{%
\subsubsection{(II) Architecture and Mathematical Derivation}\label{ii-architecture-and-mathematical-derivation}}

GAIL contains two core components that play a minimax game.

\Needspace{4\baselineskip}
\begin{enumerate}
\def\labelenumi{\arabic{enumi}.}
\tightlist
\item
\Needspace{5\baselineskip}
  Discriminator \(D_\omega\):
\end{enumerate}

\begin{itemize}
\tightlist
\item
  Role: similar to the discriminator in a GAN, it determines whether an input \((s,a)\) pair comes from the expert or the agent.
\item
  Definition: the discriminator outputs \(D(s,a)\in[0,1]\).
\item
  Closer to \(0\): the discriminator believes the data come from the expert.
\item
  Closer to \(1\): the discriminator believes the data come from the agent.
\end{itemize}

\Needspace{5\baselineskip}
The discriminator's loss function is:

\[
\mathcal{L}(\omega)
=
-
\mathbb{E}_{\pi}
\bigl[
\log D_\omega(s,a)
\bigr]
-
\mathbb{E}_{\pi_E}
\bigl[
\log(1-D_\omega(s,a))
\bigr]
\]

Explanation: the discriminator wants to maximize its ability to distinguish the two. For \((s,a)\) generated by the agent, it wants \(D(s,a)\) close to \(1\); for \((s,a)\) generated by the expert, it wants \(D(s,a)\) close to \(0\). The negative signs above convert this into a minimization problem.

\Needspace{4\baselineskip}
\begin{enumerate}
\def\labelenumi{\arabic{enumi}.}
\setcounter{enumi}{1}
\tightlist
\item
\Needspace{5\baselineskip}
  Generator or policy \(\pi\):
\end{enumerate}

\begin{itemize}
\tightlist
\item
  Role: the agent itself. Its task is to generate trajectories that ``fool'' the discriminator.
\item
  Training method: GAIL converts the discriminator's output into a reward signal for reinforcement learning.
\item
\Needspace{5\baselineskip}
  Reward function:
\end{itemize}

\[
r(s,a)
=
-
\log D_\omega(s,a)
\]

\Needspace{5\baselineskip}
Explanation:

\begin{itemize}
\tightlist
\item
  If the discriminator believes the data come from the agent, namely \(D\approx 1\), then \(\log D\approx 0\) and reward \(r\approx 0\), giving a low payoff.
\item
  If the agent successfully fools the discriminator, namely \(D\approx 0\), then \(\log D\) is a large negative number, and reward \(r=-\log D\) becomes a large positive number.
\item
  By maximizing this reward, the agent is therefore forced to generate \((s,a)\) pairs that the discriminator mistakes for expert behavior.
\end{itemize}

\Needspace{4\baselineskip}
\begin{enumerate}
\def\labelenumi{\arabic{enumi}.}
\setcounter{enumi}{2}
\tightlist
\item
\Needspace{5\baselineskip}
  Algorithm workflow and advantages:
\end{enumerate}

\begin{itemize}
\tightlist
\item
  Interactive learning: unlike BC, GAIL requires the agent to interact with the environment and sample states, then update its policy using a reinforcement learning algorithm such as PPO or TRPO.
\item
  Addresses compounding errors: GAIL matches the entire occupancy-measure distribution, rather than one-step errors. Even if the agent deviates from a trajectory, the discriminator provides feedback that encourages a return to the correct distribution. This gives GAIL long-term consistency and much greater robustness than BC.
\end{itemize}

\FloatBarrier
\Needspace{5\baselineskip}
\hypertarget{references-62}{%
\subsection{References}\label{references-62}}

\vspace{0.25em}
\begin{itemize}
\tightlist
\item
  \href{https://hrl.boyuai.com/}{Hands-on Reinforcement Learning (translated title; in Chinese)}.
\item
  Ross, S., Gordon, G. J., \& Bagnell, J. A. (2011). \href{https://proceedings.mlr.press/v15/ross11a.html}{A Reduction of Imitation Learning and Structured Prediction to No-Regret Online Learning}. AISTATS 2011.
\item
  Ho, J., \& Ermon, S. (2016). \href{https://arxiv.org/abs/1606.03476}{Generative Adversarial Imitation Learning}. NeurIPS 2016.
\item
  Silver, D., Huang, A., Maddison, C. J., et al.~(2016). \href{https://doi.org/10.1038/nature16961}{Mastering the game of Go with deep neural networks and tree search}. Nature, 529, 484--489.
\end{itemize}

\clearpage
